\daytitle{24}{交变电流、变压器与电磁应用}

\begin{learningbox}{本节导读}
上一节已经知道：线圈中的磁通量变化会产生感应电动势。本节把这个思想推进到实际电力系统：线圈在磁场中匀速转动可产生正弦交变电动势，闭合电路接入负载后形成交变电流；交变电流可通过变压器升压、降压；远距离输电要用高压降低线路损耗；最后从 LC 电磁振荡过渡到电磁波、无线通信和交直流转换。
\end{learningbox}

\textbf{本节读题路线：} 遇到交流电，先区分题目给的是“瞬时值、最大值、有效值”，再判断负载是电阻、电感还是电容；遇到变压器，先确认是否按理想变压器处理，再用匝数比定电压，由负载和功率关系定电流；遇到输电，先明确题目给的是线路输入功率还是用户获得的功率，再比较升压前后的线路电流、压降、损耗和效率。

\section*{一、学习目标}
\begin{enumerate}
  \item 能区分交变电流的瞬时值、最大值、有效值、周期和频率。
  \item 会先判断题目给的是瞬时值、最大值还是有效值，再使用 $U=\dfrac{U_m}{\sqrt2}$、$I=\dfrac{I_m}{\sqrt2}$。
  \item 能定性说明电感阻碍电流变化、电容隔直通交，以及频率对感抗、容抗的影响。
  \item 能说明理想变压器中匝数比、电压比、电流比和功率守恒之间的关系。
  \item 能用线路压降、线路损耗和输电效率解释远距离输电为什么要先升压后降压。
  \item 能说明家庭电路中火线、零线、保护接地和漏电保护的物理作用。
  \item 能用能量转换路径说明整流和逆变把交流电、直流电相互转换的基本思想。
  \item 能说明理想 LC 回路中的能量转化、振荡周期以及电磁波传播的基本图像。
\end{enumerate}
\section*{二、为什么电网采用交流}

初中电学多讨论恒定电流：电流方向不变，电压也近似不变。高中电磁感应说明，线圈转动时磁通量随时间周期性变化，感应电动势也会周期性变化。家庭用电中的“220 V、50 Hz”指的不是电压一直等于 $220\ \mathrm{V}$，而是正弦交流电的有效值为 $220\ \mathrm{V}$，频率为 $50\ \mathrm{Hz}$。

\textbf{科学发展视角：} 现代电网的核心不只是“把电送出去”，还要在发电、升压、输电、降压和用电之间控制电压、电流、频率和损耗。发电机对应电磁感应，变压器对应交变磁通和匝数比，远距离输电对应 $I^2R$ 损耗控制；把这条能量链看清，交流电就不再是一组孤立公式。

\section*{三、交流电、变压器与输电}

\begin{conceptbox}{交变电流}

大小和方向都随时间周期性变化的电流叫交变电流。最常见、最重要的是正弦交流电：
\[
u=U_m\sin(\omega t+\varphi_u),\qquad
i=I_m\sin(\omega t+\varphi_i).
\]
其中 $\varphi_u$、$\varphi_i$ 分别是电压和电流的初相位。二者的相位差记为
\[
\varphi=\varphi_u-\varphi_i.
\]
纯电阻交流电路中，电压与电流同相，$\varphi=0$；含电感或电容时，电压和电流一般不同相。图中取 $\varphi_u=0$，只画电压随时间的变化。

\begin{center}
% TikZ 图：交变电流的瞬时值、最大值和有效值
% Day 24：交变电流
\begin{tikzpicture}[scale=0.86,>=Latex]
  \draw[->,thick] (0,0) -- (7.0,0) node[right] {$t$};
  \draw[->,thick] (0,-2.0) -- (0,2.25) node[above] {$u$};
  \draw[blue,thick,samples=120,domain=0:6.28] plot(\x,{1.45*sin(180*\x/3.14159)});
  \draw[dashed,gray] (0,1.45) -- (6.35,1.45);
  \draw[dashed,gray] (0,-1.45) -- (6.35,-1.45);
  \draw[dashed,red!70] (0,1.03) -- (6.35,1.03);
  \draw[dashed,red!70] (0,-1.03) -- (6.35,-1.03);
  \node[left] at (0,1.45) {$U_m$};
  \node[left,red!70] at (0,1.03) {$U$};
  \node[left] at (0,-1.45) {$-U_m$};
  \draw[<->] (0,-1.75) -- (6.28,-1.75) node[midway,below] {$T$};
  \node[blue] at (3.2,1.9) {$u=U_m\sin\omega t$};
  \node[align=left,font=\small] at (8.45,0.55)
    {正弦交流电：\\[2pt]
     $U=\dfrac{U_m}{\sqrt2}$\\[3pt]
     $I=\dfrac{I_m}{\sqrt2}$};
\end{tikzpicture}


\figcaption{正弦交流电：瞬时值周期性变化，有效值用于比较热效应}\label{fig:ac-sine-effective}
\end{center}

\begin{itemize}
  \item $I_m$、$U_m$ 是最大值，不是平时电器标注的值。
  \item $\omega=\dfrac{2\pi}{T}=2\pi f$ 是角频率。
  \item 正弦交流电由交变电动势驱动；电流的大小和相位还由外电路负载决定。
  \item 我国家庭交流电频率 $f=50\ \mathrm{Hz}$，周期 $T=0.02\ \mathrm{s}$。
\end{itemize}

\end{conceptbox}

\begin{conceptbox}{有效值}

交流电的有效值按热效应定义：若某交流电和某直流电在同一电阻、相同时间内产生的热量相同，则这个直流电的电流、电压就是该交流电的有效值。

因此，有效值描述的不是“某一瞬间电压是多少”，而是“这个交流电让电阻发热的本领相当于多大的直流电”。电器铭牌和交流电表通常关心这个量。

对正弦交流电：
\[
\boxed{I=\frac{I_m}{\sqrt2}},\qquad
\boxed{U=\frac{U_m}{\sqrt2}}.
\]

\textbf{理解：} 有效值不是平均值。正弦交流电一个周期内电压平均值为 0，但它能让电阻发热，所以有效值不为 0。
\end{conceptbox}

\begin{conceptbox}{电感、电容对交流的作用}

电感线圈中的自感电动势阻碍电流的变化。频率越高，电流变化越快，电感对交流的阻碍越明显。理想电感的感抗为
\[
\boxed{X_L=\omega L=2\pi fL}.
\]
因此电感具有“通低频、阻高频”的特点；对稳定直流，达到稳定状态后理想电感相当于导线。

电容器极板之间没有电荷直接穿过介质，但交变电压使极板反复充、放电，外电路中可形成交变电流。理想电容的容抗为
\[
\boxed{X_C=\frac{1}{\omega C}=\frac{1}{2\pi fC}}.
\]
因此电容具有“隔直流、通交流；阻低频、通高频”的特点。

\begin{center}
% TikZ 图：电感、电容对不同频率交流的作用
% Day 24：交变电流
\begin{tikzpicture}[scale=0.88,>=Latex,line cap=round,line join=round]
  \path[use as bounding box] (-0.45,-1.95) rectangle (10.45,2.15);

  % Inductor.
  \node[font=\small\bfseries] at (2.25,1.82) {电感：通低频，阻高频};
  \draw[thick] (0.15,0.80) -- (0.95,0.80);
  \draw[thick,decorate,decoration={coil,aspect=0.5,segment length=5pt,amplitude=4pt}]
    (0.95,0.80) -- (2.55,0.80);
  \draw[thick] (2.55,0.80) -- (3.20,0.80);
  \node[font=\scriptsize] at (1.68,0.22) {自感作用阻碍电流变化};
  \node[blue!70!black,font=\small] at (1.68,-0.48)
    {$X_L=2\pi fL$};
  \draw[->,thick] (3.55,-0.68) -- (4.72,-0.68) node[right,font=\scriptsize] {$f$};
  \draw[->,thick] (3.55,-0.68) -- (3.55,0.58) node[above,font=\scriptsize] {$X_L$};
  \draw[blue!70!black,very thick] (3.70,-0.54) -- (4.50,0.36);
  \node[blue!70!black,font=\scriptsize] at (4.12,-1.22)
    {$f\uparrow\Rightarrow X_L\uparrow$};

  % Divider.
  \draw[gray!55,dashed] (5.15,-1.55) -- (5.15,1.92);

  % Capacitor.
  \node[font=\small\bfseries] at (7.72,1.82) {电容：隔直流，通交流};
  \draw[thick] (5.48,0.80) -- (6.48,0.80);
  \draw[thick] (6.48,0.27) -- (6.48,1.33);
  \draw[thick] (6.83,0.27) -- (6.83,1.33);
  \draw[thick] (6.83,0.80) -- (7.65,0.80);
  \node[font=\scriptsize] at (6.66,0.22) {极板反复充、放电};
  \node[red!75!black,font=\small] at (6.66,-0.48)
    {$X_C=\dfrac{1}{2\pi fC}$};
  \draw[->,thick] (8.62,-0.68) -- (9.82,-0.68) node[right,font=\scriptsize] {$f$};
  \draw[->,thick] (8.62,-0.68) -- (8.62,0.58) node[above,font=\scriptsize] {$X_C$};
  \draw[red!75!black,very thick]
    (8.76,0.40) .. controls (8.96,-0.18) and (9.34,-0.48) .. (9.64,-0.56);
  \node[red!75!black,font=\scriptsize] at (9.22,-1.22)
    {$f\uparrow\Rightarrow X_C\downarrow$};
\end{tikzpicture}


\figcaption{电感和电容的频率响应：感抗随频率增大，容抗随频率减小}\label{fig:ac-inductor-capacitor}
\end{center}

\textbf{功率提醒：} 交流电压和电流的有效值分别为 $U$、$I$ 时，有功功率为
\[
\boxed{P=UI\cos\varphi},
\]
其中 $\cos\varphi$ 叫功率因数。纯电阻负载中 $\varphi=0$，所以 $P=UI$；含电感、电容或电动机时，不能只把 $UI$ 当作实际消耗的平均功率。

\end{conceptbox}

\begin{conceptbox}{理想变压器}

变压器利用交变电流产生变化磁场，再由变化磁场在副线圈中产生感应电动势。在忽略铜损、铁损和漏磁的理想变压器模型中，按有效值讨论时满足
\[
\boxed{\frac{U_2}{U_1}=\frac{n_2}{n_1}},\qquad
\boxed{P_1=P_2}.
\]

\begin{center}
% TikZ 图：理想变压器模型
% Day 24：变压器
\begin{tikzpicture}[scale=0.88,>=Latex]
  \draw[fill=gray!16,draw=gray!60] (2.6,-1.7) rectangle (3.5,1.7);
  \foreach \y in {-1.1,-0.55,0,0.55,1.1}{
    \draw[blue,thick] (2.6,\y) .. controls (2.0,\y+0.25) and (2.0,\y-0.25) .. (2.6,\y-0.02);
    \draw[red,thick] (3.5,\y) .. controls (4.1,\y+0.25) and (4.1,\y-0.25) .. (3.5,\y-0.02);
  }
  \draw[->,blue,thick] (0.5,0.55) -- (1.75,0.55) node[midway,above] {$U_1$};
  \draw[->,blue,thick] (1.75,-0.55) -- (0.5,-0.55) node[midway,below] {$I_1$};
  \draw[->,red,thick] (4.35,0.55) -- (5.75,0.55) node[midway,above] {$U_2$};
  \draw[->,red,thick] (5.75,-0.55) -- (4.35,-0.55) node[midway,below] {$I_2$};
  \node[blue] at (1.55,1.55) {$n_1$};
  \node[red] at (4.55,1.55) {$n_2$};
  \node at (3.05,-2.05) {铁芯中磁通量周期性变化};
  \node[align=left,font=\small] at (8.0,0.45)
    {理想变压器：\\[2pt]
     $\dfrac{U_2}{U_1}=\dfrac{n_2}{n_1}$\\[4pt]
     $U_1I_1\approx U_2I_2$};
\end{tikzpicture}


\figcaption{理想变压器：共同铁芯让原、副线圈共享变化磁通量}\label{fig:transformer-principle}
\end{center}

\begin{itemize}
  \item $n_2>n_1$：升压变压器，副线圈电压较高；若传输同一有功功率，副线圈电流较小。
  \item $n_2<n_1$：降压变压器，副线圈电压较低；若传输同一有功功率，副线圈电流较大。
  \item 对纯电阻负载，$P_1=P_2$ 可写成 $U_1I_1=U_2I_2$，从而得到电流比与匝数比相反。
  \item 变压器不能直接改变恒定直流电压：接通瞬间可以出现暂态感应，稳定后副线圈没有持续感应电动势，原线圈还可能因电流过大而发热。
\end{itemize}

\textbf{电压比的来源：} 原、副线圈绕在同一铁芯上，理想情况下穿过每一匝线圈的磁通量变化相同。由法拉第电磁感应定律，
\[
\mathcal E_1=-n_1\frac{\mathrm d\Phi}{\mathrm dt},\qquad
\mathcal E_2=-n_2\frac{\mathrm d\Phi}{\mathrm dt}.
\]
因此同一时刻感应电动势之比等于匝数比，其有效值之比也等于匝数比。忽略线圈电阻时，端电压大小等于相应感应电动势的大小，便得到
\[
\frac{U_2}{U_1}=\frac{n_2}{n_1}.
\]
这样理解后，变压器公式不是孤立记忆，而是电磁感应定律在两个线圈上的应用。

\end{conceptbox}

\begin{samepage}
\section*{四、关系式与理想化条件}

\textbf{公式使用前检查：}
\nopagebreak[4]
\begin{enumerate}
  \item 先把题中给出的 $u,U_m,U$ 分别标成瞬时值、最大值或有效值；家用交流电 $220\ \mathrm{V}$ 指有效值。
  \item 含电感、电容时先看频率怎样改变 $X_L$、$X_C$，不要默认电压和电流同相。
  \item 理想变压器默认无漏磁、无能量损耗，输入有功功率等于输出有功功率。
  \item 判断升压或降压时，看匝数比；判断电流大小时，再结合功率关系。
  \item 输电题先说明 $P$ 是线路输入功率还是用户得到的功率，再比较线路电流、$IR_{\text{线}}$ 压降和 $I^2R_{\text{线}}$ 损耗。
  \item 家庭电路题先分清火线、零线和地线，再讨论保险、漏电保护和安全用电。
\end{enumerate}
\end{samepage}

\begin{center}
\begin{tabular}{lll}
\toprule
公式 & 物理意义 & 适用条件 \\
\midrule
$u=U_m\sin\omega t$ & 交流电瞬时电压 & 正弦交流电 \\
$U=\dfrac{U_m}{\sqrt2}$ & 电压有效值 & 正弦交流电 \\
$I=\dfrac{I_m}{\sqrt2}$ & 电流有效值 & 正弦交流电 \\
$X_L=2\pi fL$ & 感抗 & 理想电感、正弦稳态交流 \\
$X_C=\dfrac{1}{2\pi fC}$ & 容抗 & 理想电容、正弦稳态交流 \\
$P=UI$ & 平均功率 & 纯电阻或功率因数为 1 \\
$P=UI\cos\varphi$ & 交流有功功率 & $\varphi$ 为电压与电流的相位差 \\
$\dfrac{U_2}{U_1}=\dfrac{n_2}{n_1}$ & 变压器电压关系 & 理想变压器 \\
$P_1=P_2$ & 有功功率守恒 & 理想变压器 \\
$U_1I_1=U_2I_2$ & 电流关系 & 理想变压器、两侧功率因数相同 \\
$U_{\text{损}}=IR_{\text{线}}$ & 线路电压损失 & 输电线总电阻为 $R_{\text{线}}$ \\
$P_{\text{损}}=I^2R_{\text{线}}$ & 输电线热损耗 & 输电线总电阻为 $R_{\text{线}}$ \\
$T_{\mathrm{LC}}=2\pi\sqrt{LC}$ & LC 振荡周期 & 理想 LC 回路 \\
$c=\lambda f$ & 电磁波波速关系 & 真空中传播，$c\approx3.0\times10^8\ \mathrm{m/s}$ \\
\bottomrule
\end{tabular}
\end{center}

\begin{learningbox}{交流电与变压器的三层分析清单}
交流电和变压器题先分层，不要把峰值、有效值、电压比和功率关系混在一起：
\begin{enumerate}
  \item \textbf{读数层：} 先判断题目给的是瞬时值、最大值还是有效值。电表读数、铭牌电压和“$220\ \mathrm{V}$ 交流电”通常指有效值；正弦函数前的系数通常是最大值。
  \item \textbf{变压层：} 理想变压器先用 $\dfrac{U_2}{U_1}=\dfrac{n_2}{n_1}$ 定电压；升压、降压只看匝数比，不先猜电流。
  \item \textbf{负载层：} 副线圈电流由负载决定，再用功率守恒反推原线圈电流；输电题则先明确功率口径，再由 $I=P/(U\cos\varphi)$、$U_{\text{损}}=IR_{\text{线}}$ 和 $P_{\text{损}}=I^2R_{\text{线}}$ 判断压降与损耗。
\end{enumerate}
\textbf{简要归纳：} 交流题先定“哪个值”，变压器题先定“哪个电压”，功率和电流要等负载关系清楚后再算。
\end{learningbox}

\section*{五、交流电路与电磁振荡模型}

\subsection*{模型 1：交流电读数模型}

电压表、电流表在交流电路中通常读有效值。若题目给出“最大值”，要先转化为有效值；若题目给出“220 V 交流电”，默认是有效值。
\[
U_m=\sqrt2\,U,\qquad I_m=\sqrt2\,I.
\]

\textbf{读数分析清单：}
\begin{itemize}
  \item 看到 $u=U_m\sin\omega t$：这是瞬时值表达式，系数 $U_m$ 是最大值，$\omega$ 可用来求频率和周期。
  \item 看到电表读数、铭牌电压或“$220\ \mathrm{V}$ 交流电”：通常指有效值，可用于热效应、额定功率和变压器计算。
  \item 看到“一个周期内的代数平均值”：正弦交流电为 0，不能用它表示发热能力。
\end{itemize}

\textbf{认知提醒：} 交流电题先识别“给的是哪一种值”，再决定是否乘或除以 $\sqrt2$。不要看到 $220\ \mathrm{V}$ 就立刻认为它是最大值。

\subsection*{模型 2：电感、电容的频率选择}

判断电感、电容对不同频率交流信号的作用时，先写
\[
X_L=2\pi fL,\qquad X_C=\frac{1}{2\pi fC},
\]
再比较频率变化：
\[
\begin{array}{c|cc}
 & f\uparrow & f\downarrow\\
\hline
\text{电感} & X_L\uparrow,\ \text{更难通过} & X_L\downarrow,\ \text{更易通过}\\
\text{电容} & X_C\downarrow,\ \text{更易通过} & X_C\uparrow,\ \text{更难通过}
\end{array}
\]

\textbf{应用判断：} 扼流圈利用电感阻交流的作用；滤波电容常让交流成分通过旁路、使输出电压更平稳。这里的“通”“阻”表示相对容易或困难，不等于理想开关式的完全导通、完全截止。

\subsection*{模型 3：理想变压器}

先由匝数比求电压，再由功率关系求电流。对两侧功率因数相同的理想变压器，电压比跟匝数比同向，电流比与匝数比相反：
\[
\frac{I_2}{I_1}=\frac{U_1}{U_2}=\frac{n_1}{n_2}.
\]
实际解题时先看副线圈负载。若副线圈接纯电阻 $R_{\text{负载}}$，则
\[
I_2=\frac{U_2}{R_{\text{负载}}}.
\]
副线圈负载功率改变时，原线圈从电源吸收的功率也随之改变。

\begin{center}
% TikZ 图：理想变压器负载响应
% Day 24：交变电流与变压器
\begin{tikzpicture}[scale=0.88,>=Latex]
  \path[use as bounding box] (-0.35,-1.65) rectangle (8.55,1.55);

  \node[draw,rounded corners=2pt,fill=blue!6,minimum width=1.65cm,minimum height=0.85cm]
    (turns) at (0.75,0.45) {匝数比};
  \node[font=\small,align=center] at (0.75,-0.55)
    {$\dfrac{U_2}{U_1}=\dfrac{n_2}{n_1}$};

  \node[draw,rounded corners=2pt,fill=gray!10,minimum width=1.85cm,minimum height=0.85cm]
    (voltage) at (3.00,0.45) {副线圈电压};
  \node[font=\small] at (3.00,-0.55) {$U_2$};

  \node[draw,rounded corners=2pt,fill=red!6,minimum width=1.75cm,minimum height=0.85cm]
    (load) at (5.20,0.45) {负载};
  \node[font=\small] at (5.20,-0.55) {$R_{\text{负载}}$};

  \node[draw,rounded corners=2pt,fill=orange!8,minimum width=1.75cm,minimum height=0.85cm]
    (current) at (7.35,0.45) {电流};
  \node[font=\small] at (7.35,-0.55) {$I_2=U_2/R_{\text{负载}}$};

  \draw[->,thick] (turns) -- (voltage);
  \draw[->,thick] (load) -- (current);
  \draw[->,thick] (voltage.north) to[bend left=18] (current.north);

  \node[purple,font=\small,align=center] at (4.15,-1.35)
    {纯电阻负载、理想变压器中 $U_1I_1=U_2I_2$，\\
     副边电压与负载共同决定电流，负载功率需求会反馈到原线圈};
\end{tikzpicture}


\figcaption{理想变压器负载响应：匝数比定电压，负载定电流}\label{fig:transformer-load-response}
\end{center}

\subsection*{模型 4：远距离输电}

设升压站送入输电线的有功功率为 $P_{\text{入}}$，输电电压有效值为 $U$，且功率因数近似为 $1$，则
\[
I=\frac{P_{\text{入}}}{U}.
\]
若输电线总电阻为 $R_{\text{线}}$，线路电压损失、热损耗和用户端得到的功率分别为
\[
U_{\text{损}}=IR_{\text{线}},\qquad
P_{\text{损}}=I^2R_{\text{线}},\qquad
P_{\text{收}}=P_{\text{入}}-P_{\text{损}}.
\]
输电效率为
\[
\boxed{\eta=\frac{P_{\text{收}}}{P_{\text{入}}}
=1-\frac{P_{\text{损}}}{P_{\text{入}}}}.
\]
在 $P_{\text{入}}$ 和 $R_{\text{线}}$ 不变时，提高输电电压可以同时降低线路电流、压降和热损耗。若题目给的是用户必须得到的 $P_{\text{收}}$，则不能直接把它代入 $I=P/U$，还要结合 $P_{\text{入}}=P_{\text{收}}+I^2R_{\text{线}}$。

\begin{center}
% TikZ 图：远距离输电
% Day 24：交变电流
\begin{tikzpicture}[scale=0.9,>=Latex]
  \draw[rounded corners=2pt,fill=blue!8] (0,0) rectangle (1.5,1.0);
  \node at (0.75,0.5) {发电站};
  \draw[rounded corners=2pt,fill=gray!10] (2.4,-0.15) rectangle (3.5,1.15);
  \node at (2.95,0.5) {升压};
  \draw[rounded corners=2pt,fill=gray!10] (6.2,-0.15) rectangle (7.3,1.15);
  \node at (6.75,0.5) {降压};
  \draw[rounded corners=2pt,fill=red!8] (8.1,0) rectangle (9.6,1.0);
  \node at (8.85,0.5) {用户};
  \draw[->,thick] (1.5,0.5) -- (2.4,0.5);
  \draw[->,thick] (3.5,0.65) -- (6.2,0.65) node[midway,above] {高压低流};
  \draw[->,thick] (6.2,0.35) -- (3.5,0.35);
  \draw[->,thick] (7.3,0.5) -- (8.1,0.5);
  \node[font=\small,align=center] at (4.85,-0.78) {输电线损耗\\$P_{\text{损}}=I^2R$};
  \draw[red,thick,->] (4.85,0.18) -- (4.85,-0.35);
\end{tikzpicture}


\figcaption{远距离输电：升压降低线路电流、压降和热损耗，用户端再降压}\label{fig:power-transmission}
\end{center}

\subsection*{模型 5：家庭电路与安全保护}

远距离输电最终要降压到用户端。家庭电路中，火线与零线之间提供交流电压；分析一个瞬时方向时，可把工作电流路径画成“火线—用电器—零线”。保护地线通常接在金属外壳上，正常工作时不承载工作电流；一旦火线与外壳发生绝缘故障，保护地线形成低阻故障回路，使故障电流增大并促使保护装置迅速断开，同时减小人体两端可能出现的危险电压。

\begin{center}
% TikZ 图：家庭电路与安全保护
% Day 24：交变电流、变压器与电磁应用
\begin{tikzpicture}[scale=0.88,>=Latex,line cap=round,line join=round,
  every node/.style={font=\small}]
  \tikzset{
    wire/.style={thick},
    live/.style={thick,red!80!black},
    neutral/.style={thick,blue!70!black},
    earth/.style={thick,green!45!black},
    device/.style={draw,rounded corners=3pt,fill=gray!8,minimum width=1.35cm,minimum height=0.76cm},
    protector/.style={draw,rounded corners=3pt,fill=orange!10,minimum width=1.55cm,minimum height=0.82cm}
  }

  \node[draw,rounded corners=3pt,fill=blue!6,minimum width=1.35cm,minimum height=1.15cm] (grid) at (0,0.70) {用户端};
  \node[protector] (breaker) at (2.25,1.30) {空气开关};
  \node[protector] (rcd) at (4.55,0.70) {漏电保护};
  \node[device] (load) at (7.00,0.70) {用电器};

  \draw[live,->] (grid.east |- 0,1.30) -- (breaker.west);
  \draw[live,->] (breaker.east) -- (rcd.west |- 0,1.30);
  \draw[live,->] (rcd.east |- 0,1.30) -- (load.west |- 0,1.30);

  \draw[neutral,->] (load.west |- 0,0.10) -- (rcd.east |- 0,0.10);
  \draw[neutral,->] (rcd.west |- 0,0.10) -- (grid.east |- 0,0.10);

  \node[red!80!black,above] at (1.10,1.30) {火线 $L$};
  \node[blue!70!black,below] at (1.25,0.10) {零线 $N$};

  \draw[thick,fill=white] (6.35,-1.15) rectangle (7.65,-0.42);
  \node at (7.00,-0.80) {金属外壳};
  \draw[wire] (load.south) -- (7.00,-0.42);
  \draw[earth,->] (7.00,-1.15) -- (7.00,-2.05);
  \draw[earth] (6.55,-2.05) -- (7.45,-2.05);
  \draw[earth] (6.70,-2.22) -- (7.30,-2.22);
  \draw[earth] (6.86,-2.38) -- (7.14,-2.38);
  \node[green!45!black,right] at (7.13,-1.70) {保护地线 $PE$};

  \draw[red!80!black,dashed,->] (6.55,0.28) -- (6.72,-0.42);
  \draw[green!45!black,dashed,->] (6.72,-0.42) -- (7.00,-1.15);

  \node[font=\scriptsize,align=center] at (2.25,2.15) {总电流过大\\自动断开};
  \draw[orange!80!black,->] (2.25,1.78) -- (2.25,1.54);

  \node[font=\scriptsize,align=center] at (4.55,-1.08)
    {正常：$I_L=I_N$\\漏电：$I_L\ne I_N$};
  \draw[orange!80!black,->] (4.55,-0.42) -- (4.55,0.22);

  \node[font=\scriptsize,align=center] at (7.00,1.95)
    {电能转化为光、热\\或机械能};
\end{tikzpicture}


\figcaption{家庭电路保护模型：空气开关监测过电流，漏电保护器比较火线与零线电流}\label{fig:home-circuit-protection}
\end{center}

\begin{itemize}
  \item \textbf{过载保护：} 多个大功率用电器同时工作时，总电流增大，导线发热增强，保险丝或空气开关应断开电路。
  \item \textbf{保护接地：} 金属外壳接保护地线，绝缘故障时形成低阻故障回路，并配合保护装置切断电源。
  \item \textbf{漏电保护：} 火线和零线应共同穿过漏电保护器的检测环节。正常情况下两者电流大小相等；若部分电流经人体、外壳或其他路径流向大地，剩余电流不为零，漏电保护器应迅速切断电路。
\end{itemize}

\textbf{学习边界：} 本节先从电流路径和能量安全角度理解保护装置，不要求处理家庭配电施工细节；遇到真实用电问题必须由专业人员处理。

\subsection*{模型 6：LC 电磁振荡}

电容器储存电场能，线圈储存磁场能。理想 LC 回路中没有电阻损耗，电荷量 $q$、电流 $i$ 以及两种能量周期性变化：
\[
E_C=\frac12CU^2,\qquad E_L=\frac12LI^2.
\]
若初始时电容器带正向最大电荷，则一个周期内依次经历
\[
\begin{array}{c}
q=+q_m,\ i=0\\
\Downarrow\\
q=0,\ i=+I_m\\
\Downarrow\\
q=-q_m,\ i=0\\
\Downarrow\\
q=0,\ i=-I_m
\end{array}
\quad\longrightarrow\quad\text{回到初态}.
\]
振荡周期和频率为
\[
\boxed{T_{\mathrm{LC}}=2\pi\sqrt{LC}},\qquad
\boxed{f_{\mathrm{LC}}=\frac{1}{2\pi\sqrt{LC}}}.
\]

\begin{center}
% TikZ 图：理想 LC 回路一个周期内的四个典型状态
% Day 24：电磁振荡
\begin{tikzpicture}[scale=0.86,>=Latex,line cap=round,line join=round]
  \path[use as bounding box] (-0.45,-2.20) rectangle (11.15,2.15);

  \foreach \x/\top/\bottom in {
    0/{电场能最大}/{q=+q_m,\ i=0},
    2.85/{磁场能最大}/{q=0,\ i=+I_m},
    5.70/{电场能最大}/{q=-q_m,\ i=0},
    8.55/{磁场能最大}/{q=0,\ i=-I_m}
  }{
    \draw[rounded corners=3pt,fill=gray!5] (\x,0) rectangle ++(2.10,1.30);
    % Closed LC loop: capacitor in the left branch, inductor in the right.
    \draw[thick] (\x+0.50,0.98) -- (\x+1.60,0.98) -- (\x+1.60,0.82);
    \draw[thick,decorate,decoration={coil,aspect=0.45,segment length=4pt,amplitude=3pt}]
      (\x+1.60,0.82) -- (\x+1.60,0.48);
    \draw[thick] (\x+1.60,0.48) -- (\x+1.60,0.32) -- (\x+0.50,0.32);
    \draw[thick] (\x+0.50,0.98) -- (\x+0.50,0.76);
    \draw[thick] (\x+0.36,0.76) -- (\x+0.64,0.76);
    \draw[thick] (\x+0.36,0.54) -- (\x+0.64,0.54);
    \draw[thick] (\x+0.50,0.54) -- (\x+0.50,0.32);
    \node[font=\scriptsize\bfseries] at (\x+1.05,1.62) {\top};
    \node[font=\scriptsize] at (\x+1.05,-0.28) {$\bottom$};
  }

  % Charge and current cues.
  \node[red,font=\scriptsize] at (0.24,0.76) {$+$};
  \node[blue,font=\scriptsize] at (0.24,0.54) {$-$};
  \draw[orange!85!black,->,thick] (3.48,0.98) -- (4.18,0.98)
    node[midway,above,font=\scriptsize] {$i>0$};
  \node[blue,font=\scriptsize] at (5.94,0.76) {$-$};
  \node[red,font=\scriptsize] at (5.94,0.54) {$+$};
  \draw[orange!85!black,->,thick] (10.18,0.98) -- (9.48,0.98)
    node[midway,above,font=\scriptsize] {$i<0$};

  \draw[->,thick] (2.18,0.65) -- (2.72,0.65);
  \draw[->,thick] (5.03,0.65) -- (5.57,0.65);
  \draw[->,thick] (7.88,0.65) -- (8.42,0.65);
  \draw[->,thick] (10.65,-0.72) .. controls (7.40,-1.18) and (3.20,-1.18) ..
    (0.90,-0.72);
  \node[font=\scriptsize,fill=white,inner sep=1.5pt] at (5.55,-1.03)
    {继续演化，回到初态};
  \node[font=\small] at (5.55,-1.78)
    {$T_{\mathrm{LC}}=2\pi\sqrt{LC}$};
\end{tikzpicture}


\figcaption{LC 振荡：电场能和磁场能相互转化}\label{fig:lc-oscillation}
\end{center}

\textbf{真实回路：} 导线和线圈总有电阻，电磁能会逐渐转化为内能，振幅随时间减小，形成阻尼振荡。只有不断补充能量，才能维持稳定振荡。

\subsection*{模型 7：从电磁振荡到电磁波}

闭合 LC 回路中的电场和磁场主要集中在电容器与线圈附近，向外辐射并不显著。把振荡电路逐渐开放，使电场和磁场更充分地分布在空间中，并提高振荡频率，才更有利于向外辐射电磁波。高中阶段不需要推导麦克斯韦方程，但要形成一个清楚图像：\textbf{变化的电场可以产生磁场，变化的磁场可以产生电场，二者相互激发并向空间传播}。

\begin{center}
% TikZ 图：电磁波中电场和磁场的相互垂直振荡
% Day 24：交变电流、变压器与电磁应用
\begin{tikzpicture}[scale=0.88,>=Latex]
  \draw[->,thick] (0,0) -- (6.6,0) node[right] {传播方向 $c$};
  \draw[->,thick,blue] (0.6,-1.3) -- (0.6,1.4) node[above] {$E$};
  \draw[->,thick,red] (0.25,0) -- (-0.55,-0.8) node[left] {$B$};

  \draw[blue,thick,samples=120,domain=0:6.0]
    plot(\x,{0.9*sin(180*\x)});
  \draw[red,thick,samples=120,domain=0:6.0]
    plot(\x,{-0.55*sin(180*\x)-0.35});

  \foreach \x in {0.5,1.5,2.5,3.5,4.5,5.5}{
    \draw[blue!60,->] (\x,0) -- (\x,{0.9*sin(180*\x)});
    \draw[red!60,->] (\x,0) -- ++(-0.35,{-0.35-0.55*sin(180*\x)});
  }

  \node[align=center,font=\small] at (3.35,-1.65)
    {变化的电场和磁场相互激发，向外传播};
\end{tikzpicture}


\figcaption{电磁波模型：电场、磁场和传播方向两两垂直}\label{fig:electromagnetic-wave-bridge}
\end{center}

光也是电磁波。几何光学常把光近似看成“光线”，双缝干涉则显示光的波动性；这两种视角与本节的电磁波图像属于同一条知识链。

\textbf{电磁波谱：} 无线电波、微波、红外线、可见光、紫外线、X 射线和 $\gamma$ 射线都属于电磁波，只是频率和波长不同。从光子角度还可用 $E=h\nu$ 比较单个光子的能量；从波动角度看，在真空中频率越高、波长越短。对于同一障碍物或孔径尺度，波长越长，衍射通常越明显。

在真空中，所有电磁波传播速度相同：
\[
c=\lambda f\approx3.0\times10^8\ \mathrm{m/s}.
\]
因此频率越高，波长越短；频率越低，波长越长。这个关系把无线电通信、可见光和 X 射线统一到同一套波动语言中。

\begin{center}
\begin{tikzpicture}[xscale=1.05,yscale=0.95,>=Stealth]
  \draw[->,thick] (-3.55,0) -- (3.75,0) node[right,font=\small] {频率升高};
  \draw[->,thick] (3.25,-0.42) -- (-3.25,-0.42) node[left,font=\small] {波长增大};
  \foreach \x/\lab/\col in {-2.95/无线电波/blue!55,-1.82/微波/cyan!70!black,-0.72/红外/orange!80!black,0.20/可见光/red,1.12/紫外/purple,2.08/X射线/gray!70!black,3.12/$\gamma$射线/black}{
    \draw[thick,\col] (\x,-0.12) -- (\x,0.62);
  }
  \node[font=\scriptsize,align=center] at (-2.95,0.88) {无线\\电波};
  \node[font=\scriptsize,align=center] at (-1.82,1.10) {微波};
  \node[font=\scriptsize,align=center] at (-0.72,0.88) {红外};
  \node[font=\scriptsize,align=center,red] at (0.20,1.12) {可见光};
  \node[font=\scriptsize,align=center] at (1.12,0.88) {紫外};
  \node[font=\scriptsize,align=center] at (2.08,1.10) {X 射线};
  \node[font=\scriptsize,align=center] at (3.12,0.88) {$\gamma$ 射线};
  \draw[red,very thick] (-0.10,0.28) -- (0.50,0.28);
  \node[font=\scriptsize,red] at (0.20,-0.82) {人眼只看见很窄一段};
  \node[font=\scriptsize,align=center] at (-2.40,-1.25) {通信、雷达};
  \node[font=\scriptsize,align=center] at (-0.75,-1.25) {热成像};
  \node[font=\scriptsize,align=center] at (2.10,-1.25) {医学成像、材料检测};
  \node[font=\small,align=center] at (0,-1.88) {本质相同：都是真空中以光速传播的电磁波};
\end{tikzpicture}


\figcaption{电磁波谱：不同频段用途不同，但同属电磁波}\label{fig:em-spectrum}
\end{center}

\textbf{科学发展视角：} 通信、雷达、热成像、医学影像和天文观测使用的是不同频段的电磁波。“可见光”只是电磁波谱中很窄的一段，光学研究的是这一频段在传播、成像和干涉中的规律。读这类应用时，先看频率或波长属于哪个频段，再记住所有电磁波在真空中的传播速度相同。

\subsection*{模型 8：电磁波怎样承载信息}

无线通信不是把声音本身直接传播到远处，而是先把声音、图像或数据转化为电信号，再让高频电磁波携带这些信息传播。高频电磁波叫\textbf{载波}；把信息加载到载波上的过程叫\textbf{调制}；接收端把信息从载波中取出来叫\textbf{解调}。

\begin{center}
% TikZ 图：电磁波通信的基本链路
% Day 24：交变电流、变压器与电磁应用
\begin{tikzpicture}[scale=0.92,>=Latex]
  \tikzset{
    block/.style={draw,thick,rounded corners=3pt,fill=gray!7,
      minimum width=1.75cm,minimum height=0.85cm,align=center,font=\small},
    arr/.style={->,thick},
    waveblue/.style={blue,very thick},
    wavered/.style={red,very thick}
  }

  \node[block] (info) at (-5.1,0) {声音/数据\\低频信号};
  \node[block] (mod) at (-2.55,0) {调制\\装到载波上};
  \node[block] (tx) at (-0.10,0) {发射天线};
  \node[block] (rx) at (2.90,0) {接收天线};
  \node[block] (demod) at (5.45,0) {解调\\取回信息};

  \draw[arr] (info) -- (mod);
  \draw[arr] (mod) -- (tx);
  \draw[arr] (rx) -- (demod);

  \draw[wavered,domain=-5.85:-4.35,samples=70]
    plot(\x,{-1.35+0.33*sin(360*(\x+5.85))});
  \node[font=\scriptsize,red] at (-5.1,-1.92) {信息信号};

  \draw[waveblue,domain=-3.30:-1.80,samples=140]
    plot(\x,{-1.35+0.26*sin(1080*(\x+3.30))});
  \draw[wavered,domain=-3.30:-1.80,samples=70]
    plot(\x,{-1.35+0.45*sin(360*(\x+3.30))});
  \node[font=\scriptsize,align=center] at (-2.55,-2.08) {高频载波\\随信息改变};

  \draw[arr,blue] (0.72,0.72) -- (2.08,0.72);
  \draw[waveblue,domain=0.75:2.05,samples=120]
    plot(\x,{0.72+0.22*sin(900*(\x-0.75))});
  \draw[waveblue,domain=0.75:2.05,samples=120]
    plot(\x,{0.72-0.22*sin(900*(\x-0.75))});
  \draw[blue,thick] (tx.north) -- (0.55,0.58);
  \draw[blue,thick] (2.25,0.58) -- (rx.north);
  \node[font=\scriptsize,blue] at (1.40,1.28) {电磁波在空间传播};

  \draw[wavered,domain=4.70:6.20,samples=70]
    plot(\x,{-1.35+0.33*sin(360*(\x-4.70))});
  \node[font=\scriptsize,red] at (5.45,-1.92) {还原信息};

  \node[font=\small,align=center] at (0,-2.75)
    {核心思想：用高频电磁波作“载波”，把低频信息带到远处，再在接收端取出};
\end{tikzpicture}


\figcaption{电磁波通信：调制让载波携带信息，解调把信息取回}\label{fig:em-communication-chain}
\end{center}

\textbf{学习边界：} 高中阶段不要求会分析具体通信电路，但要会分清三点：载波频率高不表示波速更快；天线的典型尺寸与工作波长有关，高频载波可使用较小的天线；可传输的信息量还取决于可用带宽、调制方式、信噪比和编码，不能只由载波频率高低判断。远距离传播还受功率、天线、地形、大气和中继方式影响。

\subsection*{模型 9：交直流转换}

传统变压器只能直接处理交流电，但现代电子设备、电动汽车、光伏发电和高压直流输电经常需要在交流和直流之间转换。这里不要求会分析复杂电路，本节先抓住两句话：
\begin{itemize}
  \item \textbf{整流：} 利用二极管的单向导电性，把交流电变成方向基本不变的脉动直流。图中采用二极管整流桥得到全波整流波形，再通过滤波使电压更平稳。
  \item \textbf{逆变：} 利用开关电路按一定频率快速切换，把直流电重新变成交流电，便于接入电网或驱动交流电机。
\end{itemize}

\begin{center}
% TikZ 图：交流、直流与电力电子转换
% Day 24：交变电流、变压器与电磁应用
\begin{tikzpicture}[scale=0.88,>=Latex]
  \tikzset{
    box/.style={draw,rounded corners=3pt,thick,fill=gray!7,
      minimum width=2.35cm,minimum height=1.22cm,align=center,font=\small},
    wave/.style={very thick,blue},
    redwave/.style={very thick,red},
    arrow/.style={->,thick}
  }

  \node[box] (ac) at (0,0) {交流电\\方向周期变化};
  \node[box] (rect) at (3.4,0) {全波整流\\二极管整流桥};
  \node[box] (dc) at (6.8,0) {直流电\\方向基本不变};
  \node[box] (inv) at (10.2,0) {逆变\\直流变交流};

  \draw[arrow] (ac) -- (rect);
  \draw[arrow] (rect) -- (dc);
  \draw[arrow] (dc) -- (inv);

  \begin{scope}[shift={(0,-1.75)}]
    \draw[->] (-1.15,0) -- (1.15,0) node[right] {$t$};
    \draw[->] (-1.05,-0.58) -- (-1.05,0.72) node[above] {$u$};
    \draw[wave,domain=-1.0:1.0,samples=60]
      plot(\x,{0.45*sin(360*\x)});
  \end{scope}

  \begin{scope}[shift={(3.4,-1.75)}]
    \draw[->] (-1.15,0) -- (1.15,0) node[right] {$t$};
    \draw[->] (-1.05,-0.10) -- (-1.05,0.78) node[above] {$u$};
    \draw[redwave,domain=-1.0:1.0,samples=80]
      plot(\x,{0.45*abs(sin(360*\x))});
  \end{scope}

  \begin{scope}[shift={(6.8,-1.75)}]
    \draw[->] (-1.15,0) -- (1.15,0) node[right] {$t$};
    \draw[->] (-1.05,-0.10) -- (-1.05,0.78) node[above] {$u$};
    \draw[redwave] (-0.98,0.46) .. controls (-0.55,0.55) and (-0.20,0.37) .. (0.18,0.46)
      .. controls (0.55,0.55) and (0.82,0.39) .. (1.02,0.46);
    \draw[dashed] (-0.98,0.46) -- (1.02,0.46);
  \end{scope}

  \begin{scope}[shift={(10.2,-1.75)}]
    \draw[->] (-1.15,0) -- (1.15,0) node[right] {$t$};
    \draw[->] (-1.05,-0.58) -- (-1.05,0.72) node[above] {$u$};
    \draw[wave,domain=-1.0:1.0,samples=60]
      plot(\x,{0.45*sin(360*\x)});
  \end{scope}

  \node[font=\small,align=center] at (5.1,-3.05)
    {手机充电器、光伏并网、电动汽车和高压直流输电都需要交直流转换};
\end{tikzpicture}


\figcaption{交直流转换：整流把交流变直流，逆变把直流变交流}\label{fig:ac-dc-conversion}
\end{center}

\textbf{知识连接：} Day 19 中二极管的单向导电性，是整流的基础；本节的有效值、功率和输电损耗，则帮助理解为什么电力系统要在不同场景选择交流或直流。

\section*{六、例题精讲}

\begin{examplebox}{交流电最大值与有效值}
\textbf{题意分析：} 某正弦交流电压最大值为 $311\ \mathrm{V}$，频率为 $50\ \mathrm{Hz}$。求有效值和周期。

\textbf{物理模型：} 正弦交流电。

\begin{center}
% TikZ 图：交流电最大值与有效值例题
% Day 24：交变电流
\begin{tikzpicture}[scale=0.78,>=Latex,line cap=round,line join=round]
  \path[use as bounding box] (-0.55,-1.55) rectangle (6.65,1.80);
  \draw[->,thick] (0,0) -- (6.30,0) node[right] {$t$};
  \draw[->,thick] (0,-1.30) -- (0,1.55) node[above] {$u$};
  \draw[blue,thick,samples=120,domain=0:6.0] plot(\x,{1.10*sin(60*\x)});
  \draw[dashed,gray] (0,1.10) -- (6.0,1.10);
  \draw[dashed,red!75] (0,0.78) -- (6.0,0.78);
  \draw[dashed,gray] (0,-1.10) -- (6.0,-1.10);
  \node[left] at (0,1.10) {$311\ \mathrm{V}$};
  \node[left,red!75] at (0,0.78) {$220\ \mathrm{V}$};
  \node[left] at (0,-1.10) {$-311\ \mathrm{V}$};
  \draw[<->] (0,-1.35) -- (6.0,-1.35) node[midway,below] {$T=0.020\ \mathrm{s}$};
  \node[blue,font=\small] at (3.6,1.45) {$u=U_m\sin\omega t$};
\end{tikzpicture}

\end{center}

\textbf{解题过程：}
\[
U=\frac{U_m}{\sqrt2}=\frac{311}{1.414}\approx220\ \mathrm{V}.
\]
\[
T=\frac1f=\frac1{50}=0.020\ \mathrm{s}.
\]

\textbf{方法总结：} 家庭电路中常说的 $220\ \mathrm{V}$ 就是有效值，最大值约为 $311\ \mathrm{V}$。
\end{examplebox}

\begin{examplebox}{理想变压器}
\textbf{题意分析：} 一理想变压器原线圈接 $220\ \mathrm{V}$ 交流电，原、副线圈匝数分别为 $n_1=1100$、$n_2=110$。若副线圈接入电阻后电流为 $2.0\ \mathrm{A}$，求副线圈电压和原线圈电流。

\textbf{物理模型：} 理想降压变压器。

\begin{center}
% TikZ 图：理想变压器例题的完整电路
% Day 24：交变电流
\begin{tikzpicture}[scale=0.86,>=Latex,line cap=round,line join=round]
  % Closed core.
  \draw[gray!45,line width=9pt,rounded corners=2pt]
    (3.10,-1.08) rectangle (5.10,1.08);

  % Windings.
  \draw[blue,thick,decorate,decoration={coil,aspect=0.42,segment length=4.7pt,amplitude=4pt}]
    (2.72,-0.72) -- (2.72,0.72);
  \draw[red,thick,decorate,decoration={coil,aspect=0.42,segment length=6pt,amplitude=4pt}]
    (5.48,-0.48) -- (5.48,0.48);

  % AC source and primary connection.
  \draw[thick] (0.65,0.72) -- (2.72,0.72);
  \draw[thick] (0.65,-0.72) -- (2.72,-0.72);
  \draw[thick] (0.65,0.72) -- (0.65,0.33);
  \draw[thick] (0.65,-0.33) -- (0.65,-0.72);
  \draw[thick] (0.65,0) circle (0.33);
  \node[font=\scriptsize] at (0.65,0) {$\sim$};

  % Resistive load and secondary connection.
  \draw[thick] (5.48,0.48) -- (7.25,0.48);
  \draw[thick] (5.48,-0.48) -- (7.25,-0.48);
  \draw[thick] (7.25,0.48) -- (7.25,0.30);
  \draw[thick] (7.25,-0.30) -- (7.25,-0.48);
  \node[draw,minimum width=0.85cm,minimum height=0.72cm,font=\small] at (7.25,0) {$R$};

  % Labels, deliberately separated.
  \node[blue,font=\small] at (1.45,1.18) {$U_1=220\ \mathrm V$};
  \node[blue,font=\small] at (2.20,1.65) {$n_1=1100$};
  \node[red,font=\small] at (5.95,1.65) {$n_2=110$};
  \node[red,font=\small,anchor=west] at (7.85,0.35) {纯电阻负载};
  \node[red,font=\small,anchor=west] at (7.85,-0.25) {$I_2=2.0\ \mathrm A$};
  \node[font=\small] at (4.10,-1.52) {理想降压变压器};
\end{tikzpicture}

\end{center}

\textbf{解题过程：}
\[
\frac{U_2}{U_1}=\frac{n_2}{n_1}=\frac{110}{1100}=0.10,
\quad U_2=22\ \mathrm{V}.
\]
理想变压器功率守恒：
\[
U_1I_1=U_2I_2,
\quad I_1=\frac{U_2I_2}{U_1}
=\frac{22\times2.0}{220}=0.20\ \mathrm{A}.
\]

\textbf{方法总结：} 降压时副线圈电流可以变大，但要同时检查功率关系：纯电阻负载的理想变压器满足 $U_1I_1=U_2I_2$，电流变大来自电压降低，不是能量增加。
\end{examplebox}

\begin{examplebox}{输电损耗比较}
\textbf{题意分析：} 某升压站送入输电线的有功功率 $P_{\text{入}}=100\ \mathrm{kW}$，输电线总电阻 $R_{\text{线}}=4.0\ \Omega$，功率因数按 $1$ 处理。若输电电压分别为 $1.0\ \mathrm{kV}$ 和 $10\ \mathrm{kV}$，比较线路电流、线路压降、线路损耗和输电效率。

\textbf{物理模型：} 题目给的是线路输入功率，因此 $I=P_{\text{入}}/U$。再用
$U_{\text{损}}=IR_{\text{线}}$、$P_{\text{损}}=I^2R_{\text{线}}$ 和
$\eta=1-P_{\text{损}}/P_{\text{入}}$。

\begin{center}
% TikZ 图：不同输电电压下的压降、损耗和效率比较
% Day 24：交变电流
\begin{tikzpicture}[scale=0.82,>=Latex,line cap=round,line join=round]
  \path[use as bounding box] (-0.45,-2.15) rectangle (10.35,2.25);

  \node[draw,rounded corners=2pt,fill=blue!7,minimum width=1.25cm,minimum height=0.68cm]
    (s1) at (0.65,0.85) {输入端};
  \node[draw,rounded corners=2pt,fill=red!7,minimum width=1.25cm,minimum height=0.68cm]
    (u1) at (9.05,0.85) {接收端};
  \draw[->,thick] (s1) -- (u1);
  \node[font=\small] at (5.00,1.35)
    {$U=1.0\ \mathrm{kV},\quad I=100\ \mathrm A$};
  \node[font=\small,orange!85!black] at (3.35,0.25)
    {$U_{\text{损}}=400\ \mathrm V$};
  \node[font=\small,orange!85!black] at (6.45,0.25)
    {$P_{\text{损}}=40\ \mathrm{kW}$};
  \node[font=\small,red!75!black] at (8.65,-0.15)
    {$P_{\text{收}}=60\ \mathrm{kW},\ \eta=60\%$};

  \node[draw,rounded corners=2pt,fill=blue!7,minimum width=1.25cm,minimum height=0.68cm]
    (s2) at (0.65,-0.95) {输入端};
  \node[draw,rounded corners=2pt,fill=red!7,minimum width=1.25cm,minimum height=0.68cm]
    (u2) at (9.05,-0.95) {接收端};
  \draw[->,thick] (s2) -- (u2);
  \node[font=\small] at (5.00,-0.45)
    {$U=10\ \mathrm{kV},\quad I=10\ \mathrm A$};
  \node[font=\small,orange!85!black] at (3.35,-1.48)
    {$U_{\text{损}}=40\ \mathrm V$};
  \node[font=\small,orange!85!black] at (6.45,-1.48)
    {$P_{\text{损}}=0.40\ \mathrm{kW}$};
  \node[font=\small,red!75!black] at (8.65,-1.92)
    {$P_{\text{收}}=99.6\ \mathrm{kW},\ \eta=99.6\%$};

  \node[font=\small] at (5.15,2.00)
    {$P_{\text{入}}=100\ \mathrm{kW},\quad R_{\text{线}}=4.0\ \Omega$};
\end{tikzpicture}

\end{center}

\textbf{解题过程：}
\[
\begin{aligned}
U_1&=1.0\ \mathrm{kV}:\quad
I_1=\frac{100000}{1000}=100\ \mathrm A,\\
U_{\text{损}1}&=100\times4=400\ \mathrm V,\qquad
P_{\text{损}1}=100^2\times4=40000\ \mathrm W.
\end{aligned}
\]
\[
P_{\text{收}1}=60\ \mathrm{kW},\qquad \eta_1=60\%.
\]
\[
\begin{aligned}
U_2&=10\ \mathrm{kV}:\quad
I_2=\frac{100000}{10000}=10\ \mathrm A,\\
U_{\text{损}2}&=10\times4=40\ \mathrm V,\qquad
P_{\text{损}2}=10^2\times4=400\ \mathrm W.
\end{aligned}
\]
\[
P_{\text{收}2}=99.6\ \mathrm{kW},\qquad \eta_2=99.6\%.
\]

\textbf{方法总结：} 在线路输入功率和线路电阻不变时，电压提高 10 倍，电流和线路压降均降为原来的 $\dfrac1{10}$，线路损耗降为原来的 $\dfrac1{100}$，输电效率显著提高。先辨认功率口径，再列电流、压降、损耗和效率。
\end{examplebox}

\begin{examplebox}{LC 振荡频率与电磁波波长}
\textbf{题意分析：} 理想 LC 回路中，$L=2.5\times10^{-4}\ \mathrm H$，
$C=1.0\times10^{-10}\ \mathrm F$。求振荡频率；若以同频率向真空辐射电磁波，求波长。

\textbf{物理模型：} 先由 $L$、$C$ 决定电磁振荡频率，再用 $c=\lambda f$ 连接电路振荡与空间传播。

\textbf{解题过程：}
\[
f=\frac{1}{2\pi\sqrt{LC}}
=\frac{1}{2\pi\sqrt{2.5\times10^{-4}\times1.0\times10^{-10}}}
\approx1.0\times10^6\ \mathrm{Hz}.
\]
\[
\lambda=\frac{c}{f}
=\frac{3.0\times10^8}{1.0\times10^6}
\approx3.0\times10^2\ \mathrm m.
\]

\textbf{方法总结：} $L$、$C$ 决定振荡频率；电磁波离开电路后，在真空中仍保持同一频率，波长由 $\lambda=c/f$ 决定。
\end{examplebox}

\section*{七、变式训练}

\begin{enumerate}
  \item 某正弦交流电流有效值为 $5.0\ \mathrm{A}$，求最大值。
  \item 理想变压器原线圈 $n_1=800$ 匝，副线圈 $n_2=200$ 匝，原线圈接 $220\ \mathrm{V}$，求副线圈电压。
  \begin{center}
  \begin{tikzpicture}[scale=0.85,>=Stealth]
  \draw[thick] (0,0) -- (0.9,0);
  \draw[thick] (3.1,0) -- (4.0,0);
  \draw[thick] (1.45,-0.75) -- (1.45,0.75);
  \draw[thick] (2.55,-0.75) -- (2.55,0.75);
  \foreach \y in {-0.55,-0.33,-0.11,0.11,0.33,0.55}
    \draw[blue] (1.0,\y) arc[start angle=90,end angle=270,radius=0.11];
  \foreach \y in {-0.45,-0.15,0.15,0.45}
    \draw[red] (2.9,\y) arc[start angle=90,end angle=270,radius=0.11];
  \node[blue] at (0.55,0.42) {$U_1$};
  \node[blue,font=\small] at (0.85,1.08) {$n_1=800$};
  \node[red,font=\small] at (3.15,1.08) {$n_2=200$};
  \node[red] at (3.55,0.42) {$U_2$};
  \node at (2.0,-1.08) {$\dfrac{U_2}{U_1}=\dfrac{n_2}{n_1}$};
\end{tikzpicture}

  \end{center}
  \item 输送同样功率时，输电电压变为原来的 5 倍，输电线损耗变为原来的多少？
  \begin{center}
  \begin{tikzpicture}[scale=0.85,>=Stealth]
  \node[draw,rounded corners=2pt,fill=blue!6] (s1) at (0,0.7) {电源};
  \node[draw,rounded corners=2pt,fill=red!6] (u1) at (4.4,0.7) {用户};
  \draw[->,thick] (s1) -- (u1) node[midway,above] {$U,\ I$};
  \node[orange] at (2.2,0.25) {$P_{\text{损}}=I^2R$};
  \node[draw,rounded corners=2pt,fill=blue!6] (s2) at (0,-0.8) {电源};
  \node[draw,rounded corners=2pt,fill=red!6] (u2) at (4.4,-0.8) {用户};
  \draw[->,thick] (s2) -- (u2) node[midway,above] {$5U,\ I/5$};
  \node[orange] at (2.2,-1.25) {$P'_{\text{损}}=P_{\text{损}}/25$};
  \node[font=\small] at (2.2,1.55) {输送功率一定，升压降流};
\end{tikzpicture}

  \end{center}
  \item 家庭电路中某纯电阻用电器额定功率为 $1100\ \mathrm{W}$，接在 $220\ \mathrm{V}$ 交流电路中，正常工作电流约为多少？若多个大功率用电器同时工作，为什么要注意总电流？
  \item 一个 LC 回路中电容器电压最大时，线圈中的电流是否最大？说明能量分布。
  \item 手机充电器为什么通常先把交流电整流成直流电，再给电池充电？
  \item 某电磁波频率为 $1.0\times10^8\ \mathrm{Hz}$，在真空中的波长是多少？
  \item 保持电感 $L$ 不变，把理想 LC 回路的电容 $C$ 变为原来的 $4$ 倍，振荡周期和频率分别怎样变化？
\end{enumerate}

\begin{learningbox}{变式训练参考结果}
\begin{enumerate}[itemsep=0.45em,topsep=0.3em,parsep=0pt]
  \item $I_m=\sqrt2I=5\sqrt2\ \mathrm A\approx7.07\ \mathrm A$。
  \item $U_2/U_1=n_2/n_1$，故 $U_2=55\ \mathrm V$，为降压变压器。
  \item 在线路输入功率、功率因数和线路电阻不变时，$I=P_{\text{入}}/U$。电压变为原来的 5 倍，电流和线路压降变为原来的 $1/5$，线路损耗变为原来的 $1/25$。
  \item 正常工作电流约为 $I=P/U=5.0\ \mathrm A$。多个用电器同时工作时，总电流相加；若超过导线、开关或保护装置的额定值，发热和过载风险增大。
  \item 电容器电压最大时，电场能最大，线圈电流为零；线圈电流最大时，磁场能最大，电容器电压为零。
  \item 电池充电需要方向不变且电压、电流受控的直流电。充电器先整流，再经滤波、稳压或变换电路提供合适的直流输出。
  \item $\lambda=c/f=3.0\ \mathrm m$。
  \item 由 $T=2\pi\sqrt{LC}$，周期变为原来的 $2$ 倍；频率变为原来的 $1/2$。
\end{enumerate}
\end{learningbox}

\section*{八、有效值与电力系统拓展}
\begin{competitionbox}
\textbf{为什么有效值里会出现 $\sqrt2$？}

电阻发热由 $i^2R$ 决定。正弦电流 $i=I_m\sin\omega t$ 在一个周期内的平均发热功率为
\[
\overline P=\overline{i^2R}=I_m^2R\cdot\overline{\sin^2\omega t}.
\]
在一个周期内，
\[
\sin^2\omega t=\frac{1-\cos2\omega t}{2},
\qquad
\overline{\sin^2\omega t}=\frac12,
\]
所以
\[
\overline P=\frac12I_m^2R.
\]
若用直流有效电流 $I$ 产生同样热功率，则
\[
I^2R=\frac12I_m^2R,\qquad I=\frac{I_m}{\sqrt2}.
\]

\textbf{进一步思考：} 这个推导说明，有效值不是人为规定的换算，而是由能量和热效应决定的物理量。
\end{competitionbox}

\subsection*{1.\quad 非正弦周期电流的有效值}

有效值的热效应定义适用于任意周期电流。若周期为 $T$，则
\[
\boxed{I_{\mathrm{eff}}
=\sqrt{\frac1T\int_0^T i^2(t)\,\mathrm dt}}.
\]
对于分段恒定波形，可把一个周期分段求和：
\[
I_{\mathrm{eff}}
=\sqrt{\frac{I_1^2\Delta t_1+I_2^2\Delta t_2+\cdots}{T}}.
\]

\begin{center}
% TikZ 图：非正弦周期电流的有效值
% Day 24：竞赛入门
\begin{tikzpicture}[scale=0.92,>=Latex,line cap=round,line join=round]
  \draw[->,thick] (0,0) -- (6.65,0) node[right] {$t$};
  \draw[->,thick] (0,-1.10) -- (0,1.75) node[above] {$i$};

  \draw[blue,very thick]
    (0,1.25) -- (2.05,1.25) -- (2.05,-0.72) -- (6.15,-0.72);
  \draw[gray,dashed] (2.05,1.48) -- (2.05,-1.15);
  \draw[gray,dashed] (6.15,0.18) -- (6.15,-1.15);

  \node[left,font=\small] at (0,1.25) {$2I_0$};
  \node[left,font=\small] at (0,-0.72) {$-I_0$};
  \node[below,font=\small] at (2.05,-1.15) {$T/3$};
  \node[below,font=\small] at (6.15,-1.15) {$T$};

  \draw[orange!85!black,<->,thick] (0,1.68) -- (2.05,1.68)
    node[midway,above,font=\small] {$\Delta t_1=T/3$};
  \draw[orange!85!black,<->,thick] (2.05,-1.05) -- (6.15,-1.05)
    node[midway,below=5pt,font=\small] {$\Delta t_2=2T/3$};

  \node[font=\small,align=center] at (4.45,1.28)
    {$I_{\mathrm{eff}}^2T
      =\displaystyle\sum_k I_k^2\Delta t_k$};
\end{tikzpicture}


\figcaption{非正弦周期电流：有效值由一个周期内 $i^2$ 的平均值决定}\label{fig:ac-nonsinusoidal-rms}
\end{center}

\begin{examplebox}{分段交变电流的有效值}
\textbf{题意分析：} 一个周期内，电流在前 $T/3$ 为 $+2I_0$，后 $2T/3$ 为 $-I_0$。求有效值。

\textbf{解题过程：}
\[
I_{\mathrm{eff}}
=\sqrt{\frac{(2I_0)^2(T/3)+(-I_0)^2(2T/3)}{T}}
=\sqrt2\,I_0.
\]

\textbf{方法总结：} 求有效值时先平方再平均，电流方向的正负在平方后不会相互抵消；不能先求代数平均值。
\end{examplebox}

\subsection*{2.\quad 多副线圈变压器的功率分配}

理想变压器有多个副线圈时，每个副线圈电压仍由各自匝数决定，而原线圈输入功率等于各副线圈输出有功功率之和：
\[
\frac{U_k}{U_1}=\frac{n_k}{n_1},\qquad
P_1=P_2+P_3+\cdots.
\]

\begin{center}
% TikZ 图：多副线圈理想变压器
% Day 24：竞赛入门
\begin{tikzpicture}[scale=0.88,>=Latex,line cap=round,line join=round]
  \path[use as bounding box] (-0.55,-2.05) rectangle (9.55,2.15);

  % Closed magnetic core.
  \draw[gray!45,line width=9pt,rounded corners=2pt]
    (2.55,-1.25) rectangle (5.15,1.25);
  \draw[orange!80!black,->,thick] (3.25,1.25) -- (4.45,1.25)
    node[midway,above,font=\small] {$\Phi(t)$};

  % Primary and two secondary windings.
  \draw[blue,thick,decorate,decoration={coil,aspect=0.42,segment length=5pt,amplitude=4pt}]
    (2.15,-0.82) -- (2.15,0.82);
  \draw[red,thick,decorate,decoration={coil,aspect=0.42,segment length=5pt,amplitude=4pt}]
    (5.55,0.10) -- (5.55,0.95);
  \draw[red!65!black,thick,decorate,decoration={coil,aspect=0.42,segment length=5pt,amplitude=4pt}]
    (5.55,-0.95) -- (5.55,-0.10);

  % Primary source.
  \draw[thick] (0.65,0.82) -- (2.15,0.82);
  \draw[thick] (0.65,-0.82) -- (2.15,-0.82);
  \draw[thick] (0.65,0.82) -- (0.65,0.32);
  \draw[thick] (0.65,-0.32) -- (0.65,-0.82);
  \draw[thick] (0.65,0) circle (0.32);
  \node[font=\scriptsize] at (0.65,0) {$\sim$};
  \node[blue,font=\small] at (1.40,1.25) {$U_1,n_1$};

  % Upper load.
  \draw[thick] (5.55,0.95) -- (7.10,0.95);
  \draw[thick] (5.55,0.10) -- (7.10,0.10);
  \node[draw,minimum width=1.15cm,minimum height=0.65cm,font=\small] (Rtwo) at (7.65,0.52) {$R_2$};
  \draw[thick] (7.10,0.95) -- (Rtwo.north west);
  \draw[thick] (7.10,0.10) -- (Rtwo.south west);
  \node[red,font=\small,anchor=west] at (8.35,0.72) {$U_2,I_2$};
  \node[red,font=\small,anchor=west] at (8.35,0.28) {$n_2$};

  % Lower load.
  \draw[thick] (5.55,-0.10) -- (7.10,-0.10);
  \draw[thick] (5.55,-0.95) -- (7.10,-0.95);
  \node[draw,minimum width=1.15cm,minimum height=0.65cm,font=\small] (Rthree) at (7.65,-0.52) {$R_3$};
  \draw[thick] (7.10,-0.10) -- (Rthree.north west);
  \draw[thick] (7.10,-0.95) -- (Rthree.south west);
  \node[red!65!black,font=\small,anchor=west] at (8.35,-0.32) {$U_3,I_3$};
  \node[red!65!black,font=\small,anchor=west] at (8.35,-0.76) {$n_3$};

  \node[font=\small] at (3.85,-1.72)
    {$P_1=P_2+P_3$};
\end{tikzpicture}


\figcaption{多副线圈变压器：各路电压由匝数决定，输入功率等于各路输出功率之和}\label{fig:transformer-multiple-loads}
\end{center}

\begin{examplebox}{两个副线圈同时供电}
\textbf{题意分析：} 理想变压器原线圈 $n_1=1100$ 匝，接 $U_1=220\ \mathrm V$。两个副线圈分别为 $n_2=110$ 匝、$n_3=55$ 匝，接入纯电阻后电流分别为 $I_2=2.0\ \mathrm A$、$I_3=4.0\ \mathrm A$。求原线圈电流。

\textbf{解题过程：}
\[
U_2=22\ \mathrm V,\qquad U_3=11\ \mathrm V,
\]
\[
P_1=P_2+P_3=22\times2.0+11\times4.0=88\ \mathrm W,
\]
\[
\boxed{I_1=\frac{P_1}{U_1}=\frac{88}{220}=0.40\ \mathrm A}.
\]

\textbf{方法总结：} 多个副线圈不能分别与原线圈单独写功率相等；应先把所有输出功率相加。
\end{examplebox}

\subsection*{3.\quad 输电效率的比例规律}

若线路输入功率 $P_{\text{入}}$、线路电阻 $R_{\text{线}}$ 和功率因数均保持不变，则
\[
I=\frac{P_{\text{入}}}{U},\qquad
\frac{P_{\text{损}}}{P_{\text{入}}}
=\frac{P_{\text{入}}R_{\text{线}}}{U^2}.
\]
因此输电电压提高为原来的 $k$ 倍，线路电流与压降变为原来的 $1/k$，线路损耗及损耗率变为原来的 $1/k^2$。这一比例成立的前提是比较时 $P_{\text{入}}$、$R_{\text{线}}$ 不变。

\subsection*{4.\quad 从交流电网到电力电子}

传统交流电网便于用工频变压器升降电压；现代高压直流输电则先整流、远距离传输、再逆变并网。两者并非简单的“谁更先进”：交流系统便于组网和变压，直流线路在某些远距离、大容量或海底电缆场景中可降低特定损耗并改善控制。选择方案要综合距离、容量、换流站成本、可靠性和电网结构。

\subsection*{5.\quad 分层训练}

\begin{enumerate}
  \item \textbf{基础：} 正弦交流电 $u=220\sqrt2\sin100\pi t$，求有效值、最大值、频率和周期。
  \item \textbf{频率响应：} 频率变为原来的 $2$ 倍时，理想电感的感抗和理想电容的容抗分别怎样变化？
  \item \textbf{提高：} 理想变压器原线圈接 $220\ \mathrm{V}$，副线圈接 $22\ \Omega$ 纯电阻，匝数比 $n_1:n_2=10:1$，求副线圈电流和原线圈电流。
  \item \textbf{综合：} 在线路输入功率和线路电阻不变时，输电电压提高为原来的 $k$ 倍，线路压降、热损耗分别变为原来的多少？
  \item \textbf{波形：} 某周期电流在前半周期为 $+3I_0$、后半周期为 $-I_0$，求有效值。
  \item \textbf{科学应用：} 简述调制和解调的作用，并说明载波频率高为什么不等于传播速度大。
\end{enumerate}

\begin{learningbox}{竞赛入门训练参考结果}
\begin{enumerate}[itemsep=0.45em,topsep=0.3em,parsep=0pt]
  \item $U=220\ \mathrm V$，$U_m=220\sqrt2\ \mathrm V$，$f=50\ \mathrm{Hz}$，$T=0.020\ \mathrm s$。
  \item $X_L=2\pi fL$，故变为原来的 $2$ 倍；$X_C=1/(2\pi fC)$，故变为原来的 $1/2$。
  \item $U_2=22\ \mathrm V$，$I_2=1.0\ \mathrm A$，$I_1=U_2I_2/U_1=0.10\ \mathrm A$。
  \item 线路电流和压降均变为原来的 $1/k$，热损耗变为原来的 $1/k^2$。
  \item $I_{\mathrm{eff}}=\sqrt{[(3I_0)^2+(-I_0)^2]/2}=\sqrt5\,I_0$。
  \item 调制把信息加载到高频载波，解调从载波中取回信息；所有电磁波在真空中的传播速度均为 $c$，载波频率高只意味着真空波长更短。
\end{enumerate}
\end{learningbox}

\section*{九、易错点}
\begin{mistakebox}
\begin{enumerate}
  \item \textbf{把有效值当最大值：} 看到交流电表读数和电器铭牌，先按有效值处理；需要最大值时再乘 $\sqrt2$。
  \item \textbf{认为电压和电流总同相：} 纯电阻负载中二者同相；含电感或电容时先比较相位，平均功率应写成 $P=UI\cos\varphi$。
  \item \textbf{认为变压器能改变直流电压：} 普通变压器不能直接用于恒定直流；接通瞬间只有暂态过程，稳定后副线圈没有持续感应电动势。
  \item \textbf{只记电压比，忘记负载：} 理想变压器先用匝数比定电压，副线圈电流由负载决定，再由功率守恒反推原线圈电流。
  \item \textbf{混淆线路输入功率与用户功率：} 先写 $P_{\text{入}}=P_{\text{收}}+I^2R_{\text{线}}$，确认题目给的是哪一个功率，再计算线路电流。
  \item \textbf{把交流有效值相乘都当作有功功率：} 先判断负载是否可按纯电阻处理。纯电阻中 $P=UI$；含电感、电容或电动机时，$UI$ 可能只是视在功率，真正平均消耗功率还要看功率因数。
  \item \textbf{把保护地线当工作零线：} 先分清正常工作电流走火线和零线；地线主要在故障时保护人和设备。
  \item \textbf{认为交流和直流不能相互转换：} 遇到“交直流转换”装置时，先辨认能量流向和器件功能：交流变直流是整流，直流变交流是逆变；再按题设决定是否计入转换效率和控制条件，不把它与普通变压器的工作混为一谈。
\end{enumerate}
\end{mistakebox}

\section*{十、今日小结}
\begin{enumerate}
  \item 正弦交流电瞬时值周期性变化，交流电压表、电流表通常读有效值。
  \item 正弦交流电满足 $U=U_m/\sqrt2$、$I=I_m/\sqrt2$。
  \item 电感的感抗随频率增大，电容的容抗随频率减小；含电抗元件时电压和电流一般不同相。
  \item 理想变压器满足电压比等于匝数比，输入有功功率等于各路输出有功功率之和。
  \item 远距离输电用高压降低线路电流、压降和 $I^2R_{\text{线}}$ 损耗，从而提高效率。
  \item 家庭电路中的过载保护、保护接地和漏电保护都可以从电流路径和安全能量角度理解。
  \item LC 振荡体现电场能和磁场能的周期性转化，理想回路满足 $T=2\pi\sqrt{LC}$。
  \item 电磁波是变化电场和变化磁场相互激发并传播的结果，光是电磁波的一种。
  \item 电磁波谱把无线电波、红外线、可见光、紫外线、X 射线等统一为不同频率的电磁波，真空中满足 $c=\lambda f$。
  \item 无线通信常用高频电磁波作载波，通过调制携带信息，再由接收端解调还原信息。
  \item 整流和逆变把交流、直流连接起来，是现代电源、电动汽车、光伏并网和高压直流输电的重要基础。
\end{enumerate}

\textbf{一分钟复盘：} 先写出正弦交流的“瞬时值—最大值—有效值—周期”关系，再用
$X_L=2\pi fL$、$X_C=1/(2\pi fC)$ 判断频率变化的影响。随后画出
“发电站—升压—输电线—降压—用户”的能量链，标出变压比、线路压降、热损耗和接收功率；
最后用“电场能最大—磁场能最大—反向电场能最大—反向磁场能最大”口述 LC 回路一个周期。
复查时尤其确认：题目给的是瞬时值、最大值还是有效值，讨论的是线路输入功率还是用户得到的功率，
以及所用变压器关系是否满足理想、交流和负载条件。
